\hdhead{41}{多边形的认识——信息窗1}
\hdsection{1}{129}{分一分。（填序号。）}
\begin{scope}[scale=.6666667,draw=black,text=black]\draw[hdoutline] (168.000,241.000)--(213.000,193.000)--(299.000,241.000)--cycle;\end{scope}
\begin{scope}[scale=.6666667,draw=black,text=black]\draw[hdoutline] (423.000,237.000)--(481.500,178.500)--(540.000,237.000)--cycle;\end{scope}
\begin{scope}[scale=.6666667,draw=black,text=black]\draw[hdoutline] (666.000,236.000)--(705.000,168.450)--(744.000,236.000)--cycle;\end{scope}
\begin{scope}[scale=.6666667,draw=black,text=black]\draw[hdoutline] (186.000,392.000)--(222.000,286.000)--(258.000,392.000)--cycle;\end{scope}
\begin{scope}[scale=.6666667,draw=black,text=black]\draw[hdoutline] (434.000,390.000)--(478.000,313.790)--(522.000,390.000)--cycle;\end{scope}
\begin{scope}[scale=.6666667,draw=black,text=black]\draw[hdoutline] (649.000,315.000)--(695.000,389.000)--(754.000,389.000)--cycle;\end{scope}
\begin{scope}[scale=.6666667,draw=black,text=black]\draw[line width=.45pt] (475.843,184.157)--(481.500,189.814)--(487.157,184.157);\end{scope}
\begin{scope}[scale=.6666667,draw=black,text=black]\node[inner sep=0pt,font=\normalsize,text=black] at (227.000,264.000) {①};\end{scope}
\begin{scope}[scale=.6666667,draw=black,text=black]\node[inner sep=0pt,font=\normalsize,text=black] at (482.000,264.000) {②};\end{scope}
\begin{scope}[scale=.6666667,draw=black,text=black]\node[inner sep=0pt,font=\normalsize,text=black] at (704.000,264.000) {③};\end{scope}
\begin{scope}[scale=.6666667,draw=black,text=black]\node[inner sep=0pt,font=\normalsize,text=black] at (222.000,416.000) {④};\end{scope}
\begin{scope}[scale=.6666667,draw=black,text=black]\node[inner sep=0pt,font=\normalsize,text=black] at (478.000,415.000) {⑤};\end{scope}
\begin{scope}[scale=.6666667,draw=black,text=black]\node[inner sep=0pt,font=\normalsize,text=black] at (703.000,412.000) {⑥};\end{scope}
\q{230}{438}{200}{13}{等腰三角形}
\q{572}{438}{200}{13}{等边三角形}
\begin{scope}[scale=.6666667,draw=black,text=black]\draw[rounded corners=12,dashed,line width=.6pt] (198,478) rectangle (383,583) (537,478) rectangle (718,583);\end{scope}
\hdsection{2}{614}{在下面的点子图中分别画一个直角三角形、一个锐角三角形和一个钝角三角形。}
\begin{scope}[scale=.6666667,draw=black,text=black]\foreach \i in {0,...,40}{\foreach \j in {0,...,13}{\fill[black] ({110+\i*17.5},{705+\j*17.5}) circle (1.15);}}\end{scope}
\hdsection{3}{976}{分别画出下面三角形底边上的高。}
\begin{scope}[scale=.6666667,draw=black,text=black]\draw[hdoutline] (170.000,1041.000)--(170.000,1121.000)--(286.000,1121.000)--cycle;\end{scope}
\begin{scope}[scale=.6666667,draw=black,text=black]\draw[hdoutline] (386.000,1121.000)--(510.000,1121.000)--(447.000,1052.000)--cycle;\end{scope}
\begin{scope}[scale=.6666667,draw=black,text=black]\draw[hdoutline] (597.000,1102.000)--(754.000,1039.000)--(686.000,1102.000)--cycle;\end{scope}
\begin{scope}[scale=.6666667,draw=black,text=black]\draw[line width=.45pt] (170.000,1112.000)--(179.000,1112.000)--(179.000,1121.000);\end{scope}
\begin{scope}[scale=.6666667,draw=black,text=black]\node[inner sep=0pt,font=\normalsize,text=black] at (243.000,1066.000) {底};\end{scope}
\begin{scope}[scale=.6666667,draw=black,text=black]\node[inner sep=0pt,font=\normalsize,text=black] at (404.000,1070.000) {底};\end{scope}
\begin{scope}[scale=.6666667,draw=black,text=black]\node[inner sep=0pt,font=\normalsize,text=black] at (643.000,1121.000) {底};\end{scope}
